\newglossaryentry{CPP}{name={\CPP}, description={{\CPP} programming language}}
\newglossaryentry{OODB}{name={\OODB}, description={object-oriented database}}
\newglossaryentry{RDF}{name={\RDF}, description={resource description framework}}
\newglossaryentry{SPARQL}{name={\SPARQL}, description={{\SPARQL} protocol and {\RDF} query language}}

\chapter{Related Work}\label{ch:related}
\minitoc

\section[Pattern Matching in OOP]{Pattern Matching in Object Oriented Programming}

Besides Scala, there have been many attempts to promote pattern matching in {\OO} programs.
Because the complexity and multi-paradigm trend of {\OO} languages, a one-for-all solution is difficult.
These attempts each have their own focuses.
\Etal{Geller}, in a full-length technical report~\cite{geller2010pattern},
proposed a pattern matching framework in a purely {\OO} style for a dynamic typed language: Newspeak, one of the Smalltalk family.
Patterns are like normal objects that receive messages to perform the matching operations.
This proposal focuses on providing a framework for pattern messaging and pattern combinators, such as conjunction, disjunction, sequencing and negation,
as well as the building blocks for structural matching, such as value literals and attribute keywords.
How to match a pattern is left abstract and subject to individual object implementations.
\Etal{Widemann}~\cite{y2012paisley} similarly provided a framework in a library approach requiring generics,
that focused on the interfaces and protocols to match patterns against, extract components from and binding variables to objects.
Conjunction and disjunction are also supported through the ``both'' and ``either'' combinators.
Again, how objects and components match patterns is still abstract and therefore shallow to an object.

\Etal{Solodkyy}~\cite{solodkyy2013open} proposed an efficient functional-style pattern matching for \gls{CPP} as a library,
making heavy use of {\CPP} \emph{concepts} and template metaprogramming.
A particular pattern is expressed as a template instantiation, and the properties specific to that pattern are constructed at compile-time.
Plus the smart use of virtual table pointers as type identifiers for hashing, this gives very minimal performance overhead to the extension.
Although {\CPP} extensively supports pointers and references, but the structures are of value semantics.
This library is not specific to pointer structures.
\Etal{Ryu}~\cite{ryu2010adding} provided an language extension to Fortress, an experimental language with {\OO} perspective.
The extension depends on the \emph{getters} of an object to match its components,
and focuses on run-time type tests and multimethods.
The fashion of multimethods~\cite{muschevici2008multiple} is a strong application of pattern matching in functional languages,
and it has been extensively studied in {\OO} languages in the contrast to the single-dispatch dominance~\cite{chambers1992object,ernst1998predicate,clifton2006multijava}.

Even a bold experiment was carried out in early days by \Etal{Liu}~\cite{liu2003jmatch} in their JMatch extension to Java.
It generalizes functions and operations to be able to find the arguments back from the results, by introducing the forward and \emph{backward} modes,
including automatic iteration to search for elements in a collection.
This is more aggressive than the explicit extractors in Scala to decompose objects, with some level of operation inference.
\Etal{Emir}~\cite{emir2007matching} gave a thorough review of the commonly adopted {\OO} pattern matching techniques,
including object decomposition, visitors, type tests, typecase, case classes and extractors.
All these techniques are currently available in Scala.
\Etal{Kohn}~\cite{kohn2020dynamic} proposed a complete solution of pattern matching to the dynamically typed Python language.
It lives with the existing Python so well that it will greatly simplify the style of future Python code.
This extension has been included in the official Python 3.10~\cite{bucher2020PEP634,kohn2020PEP635},
and is thus far the most powerful algebraic pattern matching implementation in the {\OO} world.
Finally, {\OO} pattern matching is widely accepted, not only in Scala and Python,
it will appear in Java~16~\cite{bierman2020JEP394} and {\Csharp}~8~\cite{microsoft2020csharp},
the two most popular production {\OO} languages.

\section{Graph-Based Formalism}

While most research on {\OO} pattern matching all has focused on type tests and object decomposition in various abstract ways that
coexist with the {\OO} paradigm, with purposes to enhance type safety, expressiveness and conciseness,
the pointer structures of the data models in {\OO} programs are also of interest.
The semantics of a program must eventually relies on how the data objects are represented~\cite{gunter1992semantics}.
Bornat~\cite{bornat2000pointer} showed three example proofs of pointer programs in Hoare Logic, with treatment of pointer aliasing
by the principle of \emph{spatial separation}. Although the mechanism remains rather low-level, it achieved a level of local reasoning on particular data structures.
\Etal{Ke}~\cite{ke2009graph} gave a small-step operational semantics of {\OO} programs based on class, object and state graphs.
It presents a formal graph model for full featured {\OO} constructs, including stacks and method frames.
\Etal{Ke}~\cite{ke2013graph} further extended the graph model to a generic type system,
to handle recursive types, type variables and type instantiations in the notions of graph morphisms and transformations.
\Etal{Zhao}~\cite{zhao2009graph} employed the transformation of class and object graphs in the refinement of {\OO} data structures,
investigating what changes of patterns in the class structure maintained the capability of providing functionalities or services.
\Etal{Zhao}~\cite{zhao2013graph} proposed a graph-based Hoare Logic for reasoning about {\OO} programs.
The Hoare proof system consists of a set of logic rules covering most {\OO} constructs, including object creation, local variable declaration and recursive method invocation.
The soundness of the logic is given, that every specification proved by the system is valid. There wasn't a proof of completeness, although the logic is believed complete.

\section{Graph Pattern Matching}

Besides the formulation of classes and objects, graphs are widely used in the abstraction of other language notions involving relations, such as query languages.
\Etal{Francis}~\cite{francis2018cypher} described the Cypher~9 query language of the Neo4j graph database.
Cypher's core data model consists of values, property graphs and tables.
The path patterns in the query language is quite similar to the objects in {\OO} languages,
where the nodes are entity patterns like types, and the edges are relations between entities like attributes.
Formal semantics is given for the core query language based on the pattern matching.
\Etal{Tong}~\cite{tong2018mapping} proposed an approach for mapping {\OO} database models into the resource description framework (\gls{RDF}).
This work gives the formal definitions of {\OO} databases and {\RDF}, and maps the query language of {\OODB} to the \gls{SPARQL} query of {\RDF} based on graph patterns.
\Etal{Li}~\cite{li2017relaxing} defined taxonomy graphs and simulation admits the ``is-a'' relation similar to the {\OO} subtyping, but on edge labels rather than node types.
By this ``is-a'' relation, pattern matching on edges is relaxed such that a descendant edge matches an ancestor edge, and the relaxation is bounded by an inheritance distance.
Although the purpose of the relaxation was to capture more sensible matches in real-life complex data graphs,
the idea can be applied to abstract relations between objects in {\OO} programs.

\endinput